\section*{Глава 2. Проектирование веб-приложения}
\addcontentsline{toc}{section}{Глава 2. Проектирование веб-приложения}
\stepcounter{section}

\subsection{Анализ целевой аудитории и ролевая модель}

Платформа WhyGame рассчитана на два типа аудитории: \textbf{разработчиков игр},
желающих публиковать проекты без финансовых барьеров, и \textbf{игроков},
заинтересованных в удобном каталоге браузерных игр. Исходя из бизнес-логики
платформы, определены пять ролей пользователей с различными уровнями доступа.

\textbf{Гость (Guest)} --- неаутентифицированный посетитель. Имеет доступ к
каталогу, страницам игр и браузерному запуску, но не может оставлять комментарии
и оценки.

\textbf{Пользователь (User)} --- зарегистрированный участник. Может запускать
и скачивать игры, оставлять комментарии и выставлять оценки, редактировать
собственный профиль.

\textbf{Разработчик (Developer)} --- роль, запрашиваемая пользователем или
назначаемая администратором. Обладает всеми правами пользователя, а также
может загружать, редактировать и удалять собственные игры, просматривать
статистику своих публикаций в личном кабинете.

\textbf{Модератор (Moderator)} --- роль, назначаемая администратором для помощи
в проверке контента. Может просматривать все поданные на модерацию игры, одобрять
или скрывать публикации, удалять нарушающие правила комментарии. Не имеет доступа
к управлению ролями пользователей.

\textbf{Администратор (Admin)} --- привилегированная роль, назначаемая вручную.
Обладает полным доступом: включает все права модератора, а также управляет ролями
пользователей и имеет доступ ко всем разделам административной панели.

Матрица прав доступа по основным операциям системы приведена в
таблице~\ref{tab:roles}.

\begin{table}[ht!]
    \centering
    \caption{Матрица прав доступа по ролям}
    \label{tab:roles}
    \small
    \begin{tabularx}{\textwidth}{|X|c|c|c|c|c|}
        \hline
        \textbf{Операция} &
        \textbf{Гость} &
        \textbf{Польз-ль} &
        \textbf{Разраб.} &
        \textbf{Модер.} &
        \textbf{Администр.} \\ \hline
        Просмотр каталога и страниц игр     & Да & Да & Да & Да & Да \\ \hline
        Запуск игры в браузере              & Да & Да & Да & Да & Да \\ \hline
        Скачивание игры                     & Да & Да & Да & Да & Да \\ \hline
        Регистрация / вход                  & Да & --  & --  & --  & -- \\ \hline
        Комментирование игр                 & --  & Да & Да & Да & Да \\ \hline
        Выставление оценки                  & --  & Да & Да & Да & Да \\ \hline
        Редактирование профиля              & --  & Да & Да & Да & Да \\ \hline
        Загрузка и публикация игры          & --  & --  & Да & --  & Да \\ \hline
        Редактирование / удаление своей игры& --  & --  & Да & --  & Да \\ \hline
        Просмотр статистики в кабинете      & --  & --  & Да & --  & Да \\ \hline
        Просмотр очереди на модерацию       & --  & --  & --  & Да & Да \\ \hline
        Одобрение / скрытие игр             & --  & --  & --  & Да & Да \\ \hline
        Удаление комментариев               & --  & --  & --  & Да & Да \\ \hline
        Управление ролями пользователей     & --  & --  & --  & --  & Да \\ \hline
    \end{tabularx}
\end{table}

\subsection{Функциональные требования}

На основе ролевой модели выделены ключевые варианты использования системы,
представленные в виде диаграммы на рисунке~\ref{fig:use_case}.

\begin{figure}[ht!]
    \centering
    \includegraphics[width=0.92\textwidth]{use_case}
    \caption{Диаграмма вариантов использования WhyGame}
    \label{fig:use_case}
\end{figure}

Диаграмма фиксирует взаимодействие пяти акторов со следующими группами
вариантов использования:
\begin{itemize}
    \item \textbf{Просмотр и запуск} --- доступны всем акторам: каталог, страница
          игры, браузерный плеер, скачивание архива;
    \item \textbf{Управление аккаунтом} --- регистрация (только гость), вход,
          выход, редактирование профиля;
    \item \textbf{Социальная активность} --- комментарии и оценки (пользователь
          и выше);
    \item \textbf{Публикация игр} --- загрузка сборки, редактирование и удаление
          (разработчик и выше);
    \item \textbf{Модерация контента} --- просмотр очереди, одобрение / скрытие
          игр, удаление комментариев (модератор и администратор);
    \item \textbf{Администрирование} --- управление ролями пользователей
          (только администратор).
\end{itemize}

Маршруты (URL-эндпоинты) приложения, соответствующие вариантам использования,
приведены в таблице~\ref{tab:routes}.

\begin{table}[ht!]
    \centering
    \caption{Маршруты веб-приложения WhyGame}
    \label{tab:routes}
    \small
    \begin{tabularx}{\textwidth}{|p{1.4cm}|p{4.5cm}|p{2.5cm}|X|}
        \hline
        \textbf{Метод} & \textbf{URL} & \textbf{Мин. роль} & \textbf{Назначение} \\ \hline
        GET  & /                        & Гость        & Главная страница (витрина игр) \\ \hline
        GET  & /catalog                 & Гость        & Каталог с фильтрацией и поиском \\ \hline
        GET  & /game/\textlangle{}id\textrangle{}             & Гость        & Страница игры: описание, плеер, оценки, комментарии \\ \hline
        POST & /game/\textlangle{}id\textrangle{}/play        & Гость        & Фиксация запуска (инкремент \texttt{play\_count}) \\ \hline
        GET  & /game/\textlangle{}id\textrangle{}/download    & Гость        & Скачивание ZIP-архива сборки \\ \hline
        GET/POST & /register            & Гость        & Регистрация нового аккаунта \\ \hline
        GET/POST & /login               & Гость        & Вход в систему \\ \hline
        GET  & /logout                  & Пользователь & Выход из системы \\ \hline
        GET/POST & /profile/edit        & Пользователь & Редактирование своего профиля \\ \hline
        POST & /game/\textlangle{}id\textrangle{}/comment     & Пользователь & Добавление комментария \\ \hline
        POST & /game/\textlangle{}id\textrangle{}/rate        & Пользователь & Выставление оценки (1--5) \\ \hline
        GET  & /dashboard               & Разработчик  & Кабинет разработчика со статистикой \\ \hline
        GET/POST & /game/upload         & Разработчик  & Загрузка новой игры \\ \hline
        GET/POST & /game/\textlangle{}id\textrangle{}/edit    & Разработчик  & Редактирование своей игры \\ \hline
        POST & /game/\textlangle{}id\textrangle{}/delete      & Разработчик  & Удаление своей игры \\ \hline
        GET  & /admin                   & Модератор     & Панель модерации (очередь игр) \\ \hline
        POST & /admin/game/\textlangle{}id\textrangle{}/approve & Модератор  & Одобрение игры для публикации \\ \hline
        POST & /admin/game/\textlangle{}id\textrangle{}/hide  & Модератор    & Скрытие игры из каталога \\ \hline
        POST & /admin/comment/\textlangle{}id\textrangle{}/delete & Модератор & Удаление нарушающего правила комментария \\ \hline
        POST & /admin/user/\textlangle{}id\textrangle{}/role  & Администратор & Назначение роли пользователю \\ \hline
    \end{tabularx}
\end{table}

\subsection{Проектирование модели данных}

База данных платформы построена на реляционной модели и включает четыре основные
сущности: \texttt{users}, \texttt{games}, \texttt{comments} и \texttt{ratings}.
ER-диаграмма, отражающая состав полей и связи между таблицами, представлена на
рисунке~\ref{fig:er_diagram}.

\begin{figure}[ht!]
    \centering
    \includegraphics[width=\textwidth]{er_diagram}
    \caption{ER-диаграмма базы данных WhyGame}
    \label{fig:er_diagram}
\end{figure}

Связи между сущностями имеют следующий характер:
\begin{itemize}
    \item \texttt{users} $\longleftrightarrow$ \texttt{games}: один разработчик
          может иметь несколько игр (1:N, внешний ключ \texttt{developer\_id});
    \item \texttt{games} $\longleftrightarrow$ \texttt{comments}: одна игра может
          иметь несколько комментариев (1:N, внешний ключ \texttt{game\_id});
    \item \texttt{games} $\longleftrightarrow$ \texttt{ratings}: одна игра может
          иметь несколько оценок (1:N, уникальность пары \texttt{game\_id} +
          \texttt{user\_id} исключает повторное голосование).
\end{itemize}

Структура таблицы \texttt{users} представлена в таблице~\ref{tab:db_users}.

\begin{longtable}{|l|l|p{3.5cm}|p{5.5cm}|}
    \caption{Структура таблицы пользователей (\texttt{users})}
    \label{tab:db_users} \\
    \hline
    \textbf{Поле} & \textbf{Тип} & \textbf{Ограничение} & \textbf{Описание} \\ \hline
    \endfirsthead
    \tabcont{\ref{tab:db_users}}
    \textbf{Поле} & \textbf{Тип} & \textbf{Ограничение} & \textbf{Описание} \\ \hline
    \endhead
    \hline
    \endfoot
    \hline
    \endlastfoot
    id            & INTEGER   & PRIMARY KEY, AUTOINCREMENT & Уникальный идентификатор пользователя \\ \hline
    username      & VARCHAR(80)  & UNIQUE, NOT NULL        & Имя пользователя для входа и отображения \\ \hline
    email         & VARCHAR(120) & UNIQUE, NOT NULL        & Адрес электронной почты \\ \hline
    password\_hash & VARCHAR(255) & NOT NULL               & Хеш пароля (werkzeug pbkdf2) \\ \hline
    role          & VARCHAR(20)  & NOT NULL, DEFAULT 'user' & Роль: guest, user, developer, admin \\ \hline
    avatar\_path  & VARCHAR(255) & NULL                    & Путь к загруженному аватару \\ \hline
    bio           & TEXT         & NULL                    & Краткая биография разработчика \\ \hline
    created\_at   & TIMESTAMP    & DEFAULT NOW()           & Дата и время регистрации \\ \hline
\end{longtable}

Структура таблицы \texttt{games} представлена в таблице~\ref{tab:db_games}.

\begin{longtable}{|l|l|p{3.5cm}|p{5.5cm}|}
    \caption{Структура таблицы игр (\texttt{games})}
    \label{tab:db_games} \\
    \hline
    \textbf{Поле} & \textbf{Тип} & \textbf{Ограничение} & \textbf{Описание} \\ \hline
    \endfirsthead
    \tabcont{\ref{tab:db_games}}
    \textbf{Поле} & \textbf{Тип} & \textbf{Ограничение} & \textbf{Описание} \\ \hline
    \endhead
    \hline
    \endfoot
    \hline
    \endlastfoot
    id             & INTEGER      & PRIMARY KEY, AUTOINCREMENT & Уникальный идентификатор игры \\ \hline
    title          & VARCHAR(200) & NOT NULL                   & Название игры \\ \hline
    description    & TEXT         & NULL                       & Описание игры \\ \hline
    developer\_id  & INTEGER      & FK → users.id, NOT NULL    & Ссылка на автора-разработчика \\ \hline
    genre          & VARCHAR(50)  & NULL                       & Жанр игры (action, puzzle и т.\,д.) \\ \hline
    cover\_path    & VARCHAR(255) & NULL                       & Путь к обложке игры \\ \hline
    build\_path    & VARCHAR(512) & NOT NULL                   & Путь к директории Unity WebGL сборки \\ \hline
    is\_downloadable & BOOLEAN    & DEFAULT FALSE              & Флаг доступности ZIP-архива \\ \hline
    status         & VARCHAR(20)  & DEFAULT 'pending'          & Статус: pending, published, hidden \\ \hline
    play\_count    & INTEGER      & DEFAULT 0                  & Счётчик запусков игры \\ \hline
    download\_count & INTEGER     & DEFAULT 0                  & Счётчик скачиваний \\ \hline
    created\_at    & TIMESTAMP    & DEFAULT NOW()              & Дата публикации \\ \hline
    updated\_at    & TIMESTAMP    & DEFAULT NOW()              & Дата последнего редактирования \\ \hline
\end{longtable}

Структура таблицы \texttt{comments} представлена в таблице~\ref{tab:db_comments}.

\begin{table}[ht!]
    \centering
    \caption{Структура таблицы комментариев (\texttt{comments})}
    \label{tab:db_comments}
    \small
    \begin{tabularx}{\textwidth}{|l|l|p{3.5cm}|X|}
        \hline
        \textbf{Поле} & \textbf{Тип} & \textbf{Ограничение} & \textbf{Описание} \\ \hline
        id         & INTEGER   & PRIMARY KEY, AUTOINCREMENT & Уникальный идентификатор комментария \\ \hline
        game\_id   & INTEGER   & FK → games.id, NOT NULL    & Игра, к которой относится комментарий \\ \hline
        author\_id & INTEGER   & FK → users.id, NOT NULL    & Пользователь, оставивший комментарий \\ \hline
        body       & TEXT      & NOT NULL                   & Текст комментария \\ \hline
        created\_at & TIMESTAMP & DEFAULT NOW()             & Дата и время публикации \\ \hline
    \end{tabularx}
\end{table}

Структура таблицы \texttt{ratings} представлена в таблице~\ref{tab:db_ratings}.

\begin{table}[ht!]
    \centering
    \caption{Структура таблицы оценок (\texttt{ratings})}
    \label{tab:db_ratings}
    \small
    \begin{tabularx}{\textwidth}{|l|l|p{3.5cm}|X|}
        \hline
        \textbf{Поле} & \textbf{Тип} & \textbf{Ограничение} & \textbf{Описание} \\ \hline
        id         & INTEGER    & PRIMARY KEY, AUTOINCREMENT        & Уникальный идентификатор оценки \\ \hline
        game\_id   & INTEGER    & FK → games.id, NOT NULL           & Оцениваемая игра \\ \hline
        user\_id   & INTEGER    & FK → users.id, NOT NULL           & Пользователь, поставивший оценку \\ \hline
        score      & SMALLINT   & CHECK(1 $\leq$ score $\leq$ 5), NOT NULL & Числовая оценка от 1 до 5 \\ \hline
        created\_at & TIMESTAMP & DEFAULT NOW()                     & Дата и время выставления оценки \\ \hline
        ---        & ---        & UNIQUE(game\_id, user\_id)        & Ограничение: не более одной оценки от пользователя за игру \\ \hline
    \end{tabularx}
\end{table}

Агрегированный рейтинг игры вычисляется как среднее арифметическое по всем
оценкам:

\begin{equation}
    \bar{R} = \frac{1}{N} \sum_{i=1}^{N} s_i,
    \label{eq:avg_rating}
\end{equation}

\begin{where}
    \item $s_i$~--- оценка $i$-го пользователя (целое число от 1 до 5);
    \item $N$~--- общее количество оценок игры.
\end{where}

\subsection{Архитектура приложения}

Приложение построено по трёхуровневой архитектуре (рисунок~\ref{fig:architecture}):
клиентский уровень (браузер), серверный уровень (Flask-приложение) и уровень
хранения данных (PostgreSQL и файловая система).

\begin{figure}[ht!]
    \centering
    \includegraphics[width=0.88\textwidth]{app_architecture}
    \caption{Трёхуровневая архитектура приложения WhyGame}
    \label{fig:architecture}
\end{figure}

Серверный уровень организован по шаблону \textbf{Application Factory}: функция
\texttt{create\_app()} создаёт и конфигурирует экземпляр приложения, подключает
расширения (SQLAlchemy, Migrate, Login Manager) и регистрирует блупринты (Blueprint).
Такой подход упрощает тестирование и позволяет создавать несколько экземпляров
приложения с разной конфигурацией.

Маршруты сгруппированы в четыре блупринта по функциональному назначению:
\textbf{auth} (регистрация, вход, выход), \textbf{games} (каталог, просмотр, загрузка,
редактирование), \textbf{users} (профили) и \textbf{admin} (панель модерации).
Каждый блупринт выделен в отдельный модуль.

Структура каталога проекта показана в листинге~\ref{lst:structure}.

\begin{lstlisting}[caption={Структура проекта WhyGame}, label={lst:structure}]
WhyGame/
  app/
    __init__.py          # Application Factory, инициализация расширений
    models.py            # SQLAlchemy-модели: User, Game, Comment, Rating
    blueprints/
      auth.py            # Маршруты: /register, /login, /logout
      games.py           # Маршруты: /catalog, /game/<id>, /game/upload
      users.py           # Маршруты: /profile, /user/<username>
      admin.py           # Маршруты: /admin, /admin/game, /admin/user
    templates/
      base.html          # Базовый шаблон (навигация, подключение CSS/JS)
      catalog.html       # Каталог игр
      game.html          # Страница игры (плеер, комментарии, оценки)
      upload.html        # Форма загрузки игры
      dashboard.html     # Личный кабинет разработчика
    static/
      css/               # Стили
      js/                # Клиентские скрипты
      uploads/           # Директория хранения игровых сборок
  config.py              # Конфигурация: SECRET_KEY, DATABASE_URL
  run.py                 # Точка запуска
  requirements.txt       # Зависимости Python
\end{lstlisting}

Хранение игровых сборок Unity WebGL организовано в файловой системе сервера.
При загрузке ZIP-архива его содержимое распаковывается в директорию
\texttt{static/uploads/<game\_id>/}. Для запуска игры Flask отдаёт входной
HTML-файл сборки через соответствующий маршрут; браузер самостоятельно загружает
файлы \texttt{.wasm} и \texttt{.data}, которые обслуживаются как статические
ресурсы.

\subsection{Проектирование интерфейса}

Интерфейс платформы проектировался из расчёта трёх ключевых сценариев использования:
просмотр каталога, запуск игры и загрузка новой публикации.

\textbf{Каталог игр} --- основная точка входа для игроков. Страница содержит
сетку карточек игр с обложкой, названием, жанром и средним рейтингом. Предусмотрены
поле поиска по названию и фильтры по жанру и статусу. В верхней части страницы
расположена навигационная панель с логотипом платформы, ссылками на разделы
и кнопками «Войти» / «Профиль». Кнопка «Загрузить игру» отображается только
пользователям с ролью разработчика или администратора.

\textbf{Страница игры} --- центральный элемент взаимодействия. Включает заголовок
и имя разработчика, встроенный \texttt{<canvas>}-плеер для Unity WebGL, блок
статистики (количество запусков и оценок), форму выставления оценки (5 звёздочек,
доступна авторизованным пользователям), текстовое описание игры, кнопку скачивания
ZIP-архива (если активирована разработчиком) и секцию комментариев с формой
добавления нового.

\textbf{Форма загрузки игры} --- доступна пользователям с ролью разработчика.
Содержит следующие поля: название игры (однострочный текст), описание
(многострочное поле), жанр (выпадающий список), загрузка обложки (JPG / PNG),
загрузка сборки Unity WebGL (ZIP-архив) и флаг «Разрешить скачивание». После
отправки формы игра переходит в статус \texttt{pending} и ожидает одобрения
модератора.

\textbf{Личный кабинет разработчика} --- доступен по адресу \texttt{/dashboard}.
Отображает список собственных игр с их статусами, счётчиками запусков, скачиваний
и средним рейтингом. Предоставляет быстрый доступ к редактированию и удалению
каждой публикации.

\subsection{Макет страницы}

Для визуализации ключевых экранов платформы разработаны макеты двух основных
страниц: каталога игр и страницы отдельной игры. Макеты отражают расположение
элементов интерфейса до начала вёрстки и служат ориентиром при реализации
HTML"=шаблонов.

\textbf{Макет каталога игр} (рисунок~\ref{fig:mockup_catalog}) демонстрирует
структуру главного экрана платформы: навигационная панель с логотипом и кнопками
авторизации, строка поиска и фильтры по жанру, а также сетка карточек игр
с обложкой, названием и средним рейтингом.

\begin{figure}[ht!]
    \centering
    \includegraphics[width=0.9\textwidth]{mockup_catalog}
    \caption{Макет страницы каталога игр WhyGame}
    \label{fig:mockup_catalog}
\end{figure}

\textbf{Макет страницы игры} (рисунок~\ref{fig:mockup_game}) отражает компоновку
страницы отдельной игры: заголовок с именем разработчика, встроенный
\texttt{<canvas>}"=плеер Unity WebGL, блок статистики (запуски, оценка),
форма выставления оценки, текстовое описание, кнопка скачивания ZIP"=архива
и секция комментариев.

\begin{figure}[ht!]
    \centering
    \includegraphics[width=0.9\textwidth]{mockup_game}
    \caption{Макет страницы игры WhyGame}
    \label{fig:mockup_game}
\end{figure}
